\documentclass[addpoints]{exam}
% \documentclass[addpoints,12pt]{exam}

\usepackage[slovene]{babel}
\usepackage{amsfonts}
\usepackage{amssymb}
\usepackage{amsmath}
\usepackage{float}
\usepackage{graphicx}
\usepackage{enumitem}
\usepackage{color}
\usepackage{eurosym}

\usepackage{pgfpages}
% \usepackage{tikz}
\usepackage{wrapfig}
\usepackage{pgfkeys}
\usepackage{pgfplots}
\usepackage{xcolor}
\usepackage{tkz-euclide}
\usepackage{xfp}




%Komentar naslednje vrstice glede na verzijo testa -- rešitve ali prostor za odgovore:
% \printanswers

% \linespread{2}


\pointsinrightmargin
\marginbonuspointname{}


\renewcommand{\solutiontitle}{\noindent\textbf{Rešitve in točkovanje:}\par\noindent}

\colorgrids
\definecolor{GridColor}{gray}{0.7}
\setlength{\gridsize}{7mm}

\extrawidth{5mm}
\setlength{\rightpointsmargin}{1.5cm}

\begin{document}

\runningfooter{}{}{}

\begin{center}
    \fbox{\fbox{\parbox{15cm}{\centering
    Popravljanje in pridobivanje ocen -- četrto pisno ocenjevanje znanja \\ 2. letnik -- test C \\ 15. junij 2026 \\ (Kvadratna funkcija)}}}
\end{center}

\vspace{0.1in}
\makebox[\textwidth]{Ime in priimek, oddelek:\enspace\hrulefill}


\begin{center}
    \hqword{Naloga}
    \hpword{Možne točke}
    \chbpword{Dodatne točke}
    \hsword{Dosežene točke}
    \htword{Skupaj}  
    % \combinedpointtable[h][questions]
    \gradetable[h][questions]

\end{center}

\vspace{0.1in}


\begin{questions}

    % 1. vprašanje

    \question
    Dana je kvadratna funkcija $f(x)=\frac{1}{4}x^2-\frac{1}{4}x-\frac{15}{16}$.

    \begin{parts}
        
        \part[4]
        Zapišite predpis funkcije $f$ v ničelni obliki.

            \begin{solution}[8cm]

            \end{solution}

            
        \part[5]
        Zapišite predpis funkcije $f$ v temenski obliki in določite teme.

            \begin{solution}[5cm]

            \end{solution}


        \newpage
        \part[1]
        Določite začetno vrednost funkcije $f$.

            \begin{solution}[3cm]

            \end{solution}


        \part[2]
        Določite enačbo simetrijske osi grafa funkcije $f$.

            \begin{solution}[3cm]

            \end{solution}


        \part[5]
        Narišite graf funkcije $f$.

            \begin{solution}[5cm]

            \end{solution}
        

        \newpage
        \part[2]
        Določite definicijsko območje in zalogo vrednosti funkcije $f$.

            \begin{solution}[3cm]

            \end{solution}
        
        \part[4]
        Graf funkcije $f$ togo premaknemo za vektor $\vec{v}=(-0.5,1)$ in tako dobimo graf funkcije $g$.
        Določite predpis funkcije $g$.

            \begin{solution}[8cm]

            \end{solution}


        \part[3]
        Določite predpis funkcije $h$, katere graf dobimo, če graf funkcije $f$ raztegnemo za faktor $3$ vzdolž ordinatne osi ter zrcalimo preko ordinatne osi.

            \begin{solution}[7cm]

            \end{solution}


    \end{parts}


    \newpage
    % 2. vprašanje

    \question
    Dana je družina parabol $y=2x^2+(3-m)x+2$.

    \begin{parts}
        
        \part[6]
        Določite vrednost parametra $m$ tako, da bo imela parabola z abscisno osjo eno skupno točko.
        Določite tudi točko dotikališča.

            \begin{solution}[13cm]

            \end{solution}


        \part[4]
        Določite $m$ tako, da parabola z abscisno osjo ne bo imela skupnih točk.

            \begin{solution}[6cm]

            \end{solution}


        % \part[2]
        % Za katero vrednost parametra $m$ enačba predstavlja parabolo? 
        % Določite presečišče parabole z ordinatno osjo.

        %     \begin{solution}[2cm]

        %     \end{solution}


    \end{parts}



    






\end{questions}



\end{document}